\SourceLeaf{283}%
\begin{center}
\begin{minipage}[t]{5.7cm}\centering $\begin{array}{rcc}
&&a\\
1.8&+&\\
&&b\\
2.9&-&\\
&&a''\\
3.18&-&\\
&&c\\
4.17&+&\\
&&a^{\mathrm{VI}}\\
5.12&+&\\
&&b^{\mathrm{V}}\\
6.21&+&\\
&&d\\
7.14&-&\\
&&b^{\mathrm{VIII}}\\
8.1&-&\\
&&a'\\
9.2&+&\\
&&b'\\
10.19&-&\\
&&a^{\mathrm{IV}}\\
11.16&+&\\
\end{array}$\end{minipage}
\hfill
\begin{minipage}[t]{5.7cm}\centering $\begin{array}{rcc}
&&b^{\mathrm{IV}}\\
12.5&-&\\
&&a^{\mathrm{VII}}\\
13.22&+&\\
&&b^{\mathrm{VII}}\\
14.7&+&\\
&&d'\\
15.20&-&\\
&&b'''\\
16.11&-&\\
&&a^{\mathrm{V}}\\
17.4&-&\\
&&c'\\
18.3&+&\\
&&a'''\\
19.10&+&\\
&&b''\\
20.15&+&\\
&&d''\\
21.6&-&\\
&&b^{\mathrm{VI}}\\
22.13&-&\\
\end{array}$\end{minipage}
\end{center}
\[
(a) = \,.\, -b+c-d
\]
\[
(b) = +a\,.\, +c-d
\]
\[
(c) = -a-b\,.\, -d
\]
\[
(d) = +a+b+c\,.
\]

\Needspace{20\baselineskip}\begin{center}[6.]\end{center}

Der Beweis der oben mitgetheilten Regel beruhet auf einer andern Vertheilung
der Figur in mehrere geschlossene Figuren, von denen keine sich selbst
schneidet, obwohl eine die andere schneiden kann; diess sind nemlich
diejenigen, die sich nach und nach bilden, während die Linie erst beschrieben
wird. Man hebe von den Intersectionen bloss diejenige Hälfte aus, wo sie
zum zweitenmale geschehen, in unserm Beispiel
